\chapter{Forsking gjennom design er kunnskap}

\begin{motto}
Eit krav om at andre skal kunne byggje vidare på arbeidet,\\ er strengare enn dei fleste vitskapar stiller til seg sjølve.
\end{motto}

Sett at ein gjev seg over på at design har metode. Eit hardare spørsmål står att: produserer design \emph{kunnskap}? Eller berre ting — vakre, nyttige, smarte ting, men ting, ikkje innsikt som andre kan etterprøve og leggje til sitt eige? Skuldinga frå den kritiske sida er at forsking gjennom design ikkje er etterprøvbar og ikkje kumulativ — at kvart prosjekt er ein eingongshending, eit kunstverk forkledd som forsking, utan noko andre kan ta med seg. Dette kapitlet svarar at det motsette er tilfellet, og at svaret ikkje er ei forsvarstale, men eit veksande forskingsfelt med eksplisitte rigorkriterium. Forsking gjennom design — det å lage noko som ein måte å vite på — er ei etablert kunnskapsform med si eiga litteratur, sine eigne standardar, og, viktigast av alt, eit innebygd krav om at kunnskapen skal kunne byggjast vidare på.

\section{Tre måtar å forske på}

Klargjeringa byrjar med eit skilje som Christopher Frayling gjorde i ein artikkel som er blitt eit referansepunkt for heile feltet \autocite{frayling1993}. Frayling skilde mellom forsking \emph{om} design — historie, teori, estetikk, der designet er studieobjektet; forsking \emph{for} design — der ein samlar kunnskap som skal nyttast i ein komande designprosess; og forsking \emph{gjennom} design, der sjølve det å designe er framgangsmåten ein vinn kunnskap på. Det siste er det omstridde. Men Fraylings poeng var nett at det ikkje er ein selvmotseiing: ein kan lære noko sant og overførbart \emph{ved} å lage, slik ein kjemikar lærer ved å eksperimentere. Forma på kunnskapen er annleis — han ligg ofte i artefaktet og det som blei oppdaga undervegs, ikkje i ein proposisjon — men det gjer han ikkje mindre verkeleg. Skiljet gav feltet eit vokabular til å seie kva det heldt på med, og det er sidan blitt utbygd av andre. \textcite{fallman2003design} skilde til dømes mellom designorientert forsking, der designet tener kunnskapsmålet, og forskingsorientert design, der kunnskapen tener designmålet — eit grep som rydda opp i ein vedvarande forvirring om kva slags påstand eit designprosjekt eigentleg gjer.

\section{Rigorkriteria}

Det avgjerande svaret på skuldinga om manglande etterprøvbarheit kom frå John Zimmerman, Jodi Forlizzi og Shelley Evenson, i ein artikkel som sette ein standard feltet sidan har målt seg mot \autocite{zimmerman2007rtd}. Dei spurde rett ut: kva skil god forsking gjennom design frå berre god design? Og dei svara med fire eksplisitte kriterium. Det fyrste er \emph{prosess} — at framgangsmåten er dokumentert og rasjonell nok til at ein annan forskar kan forstå og vurdere kvifor det blei gjort slik det blei. Det andre er \emph{oppfinning} — at arbeidet yter eit reelt, nytt bidrag, ikkje berre ein variasjon av det kjende. Det tredje er \emph{relevans} — at det adresserer noko som faktisk betyr noko. Og det fjerde, det som råkar skuldinga hardast, er \emph{utvidbarheit} (\textit{extensibility}): kravet om at kunnskapen frå prosjektet skal dokumenterast og formidlast så grundig at andre kan byggje vidare på han.

Stå ved det siste kriteriet eit augneblink, for det er meir radikalt enn det ser ut. Utvidbarheit krev kumulasjon. Eit prosjekt som oppfyller kriteriet, \emph{må} etterlate seg noko etterprøvbart og vidareførbart — elles er det per definisjon ikkje god forsking gjennom design. Med andre ord har feltet sjølv, lenge før kritikken kom, gjort kumulativitet til eit krav for kvalitet. Skuldinga om at forsking gjennom design ikkje er kumulativ, råkar då ikkje feltet; ho råkar berre arbeid som feltet sjølv ville rekna som mislukka forsking. Det er ein viktig skilnad. Ein dømer ikkje fysikken etter den dårlege fysikken, eller medisinen etter kvakksalveriet. Ein dømer eit fag etter standarden det set for godt arbeid — og her er standarden eksplisitt at andre skal kunne byggje vidare.

\section{Kva slags kunnskap det er}

Men kva \emph{slags} kunnskap produserer forsking gjennom design eigentleg? Her er William Gaver sin analyse den mest nøkterne og difor den mest overtydande \autocite{gaver2012what}. Gaver hevdar ikkje for mykje. Han innrømmer at teorien som kjem ut av forsking gjennom design er provisorisk — han er ikkje universell, ikkje endeleg, ofte bunden til dei artefakta og situasjonane han oppstod i. Men, seier Gaver, det gjer han ikkje uvitskapleg; det gjer han \emph{annorleis} vitskapleg. Teori i dette feltet er koherent og rettleiande utan å vere fullstendig eller falsifiserbar i Popper si strenge tyding — og det same, minner Gaver oss om, gjeld store delar av teorien i andre praktiske og menneskelege fag. Det avgjerande poenget hans er at vi ikkje bør krevje at forsking gjennom design skal produsere den same typen kunnskap som naturvitskapen, og så døme henne mislukka når ho ikkje gjer det. Vi bør spørje kva ho faktisk leverer — provisorisk, men koherent teori, robust nok til å rettleie vidare arbeid — og verdsetje det på sine eigne premissar. Dette er den same korreksjonen vi har sett før, men no anvendt på sjølve kunnskapsutbyttet.

At kunnskapen er overførbar og ikkje berre lokal, har vore vist konkret. Joep Frens sin doktoravhandling om rik interaksjon er eit lærebokdøme \autocite{frens2006rich}. Frens designa ein serie konsept for kameragrensesnitt — ikkje for å lage eit produkt, men for å undersøkje eit prinsipp om korleis form, interaksjon og funksjon kan integrerast. Det avhandlinga etterlét seg, var ikkje kamera; det var overførbar innsikt om eit designprinsipp, demonstrert gjennom artefakt og formulert så andre kunne nytte det i heilt andre samanhengar. Konseptdesignet var middelet; den generaliserbare kunnskapen var målet. Dette er forsking gjennom design som gjer akkurat det skuldinga seier ho ikkje kan: produserer noko andre kan ta med seg.

\section{Den nordiske konsolideringa}

Det som gjer forsvaret sterkast, er likevel ikkje dei einskilde studiane, men at feltet er blitt \emph{konsolidert} — at det no har lærebøker, undervist metode og teoretisk sjølvforståing. Tre verk, mange av dei nordiske, ber denne konsolideringa. Det fyrste er \textit{Design Research through Practice} av Ilpo Koskinen og medforfattarane hans, som samlar feltet under nemninga \emph{constructive design research} og gjev det ein klar typologi \autocite{koskinen2011constructive}: laboratoriet, der ein studerer under kontrollerte vilkår; feltet, der ein studerer i bruk og kontekst; og utstillingsrommet, der det designa artefaktet sjølv ber argumentet, slik eit kritisk eller spekulativt prosjekt gjer. Tre tilnærmingar, kvar med sine kriterium for kva som tel som godt arbeid. Dette er ikkje eit fag som famlar; det er eit fag som har kartlagt sine eigne metodar grundig nok til å undervise dei. At forsking gjennom design no \emph{er} undervist metode, dokumentert mellom anna av \textcite{stappers2017rtd} i eit standardverk for faget, er i seg sjølv eit motbevis mot at feltet er for ustøtt til å vidareførast. Ein lærer ikkje vidare det som ikkje finst.

Dei to andre verka går djupare i grunnlaget. Thomas Binder og Johan Redström la fram ideen om \emph{exemplary design research}, der kunnskapen ber gjennom samspelet mellom eit \emph{program} — eit eksplisitt formulert sett av antakingar og ambisjonar — og dei \emph{eksperimenta} som prøver programmet ut \autocite{binder2006exemplary}. Dette er ein presis mekanisme for korleis design kan akkumulere: programmet gjev det einskilde eksperimentet meining utover seg sjølv, og eksperimenta talar tilbake til programmet, justerer det, gjer det skarpare. Det er, om ein vil, korleis eit designforskingsfelt nærmar seg noko som liknar eit paradigme — ikkje ved å vente på éi stor teori, men ved å byggje koherente program og prøve dei systematisk. Og Redström førte dette heilt fram til spørsmålet om teori i \textit{Making Design Theory} \autocite{redstrom2017making}. Hans tese er at design ikkje treng å låne teorien sin frå andre fag, men \emph{lagar} sin eigen — at teoriproduksjon sjølv er ei designhandling, der omgrep og rammeverk vert forma, prøvde og revidert slik artefakt vert det. Med Redström har feltet teke det siste steget: frå å forsvare at det produserer kunnskap, til å gjere greie for korleis det produserer si eiga teori.

\vignett

Eit fag som har metode og som produserer overførbar kunnskap, bør òg kunne vise at kunnskapen faktisk hopar seg opp over tid. Det er den siste og mest konkrete prøva, og ho er emnet for neste kapittel.

\vidarelesing{\fullcite{frayling1993}; \fullcite{zimmerman2007rtd}; \fullcite{gaver2012what}; \fullcite{koskinen2011constructive}; \fullcite{binder2006exemplary}; \fullcite{redstrom2017making}.}
